%==============================================================================
% Digitalizing Port Processes: Design and Causal Evaluation of a Port
% Transformation Observatory
%
% Periódico-alvo: Government Information Quarterly (Elsevier) — decisão 40, 12/09/2026.
% Formato conforme o Guide for Authors da GIQ:
%   - classe elsarticle (LaTeX aceito), citações autor-ano (APA 7 na produção);
%   - revisão DUPLAMENTE ANÔNIMA: este arquivo NÃO contém autor/afiliação/
%     agradecimentos — tudo isso fica em title_page.tex (arquivo separado);
%   - abstract <= 250 palavras; 1–7 keywords; highlights 3–5 itens <= 85 caracteres
%     (também em arquivo separado: highlights.tex);
%   - seções numeradas; apêndices A, B, C com Eq. (A.1) etc.
%
% REGRA DO PROJETO: nenhum resultado é digitado à mão. Todo número vem de
% % numbers.tex — GERADO AUTOMATICAMENTE por scripts/12_export_overleaf.R.
% Não editar à mão: qualquer alteração é sobrescrita na próxima exportação.
% Gerado em 2026-09-23 12:08:40.72706 a partir de data/processed/observatory.duckdb

\newcommand{\Natracacoes}{1,182,270}
\newcommand{\Anoinicio}{2010}
\newcommand{\Anofim}{2026}
\newcommand{\Npainellinhas}{29,461}
\newcommand{\Nportos}{251}
\newcommand{\Nportospublicos}{34}
\newcommand{\Nportostups}{217}
\newcommand{\Fretemedio}{4.9\%}
\newcommand{\Fretepico}{7.2\%}
\newcommand{\Anofretepico}{2022}
\newcommand{\Nfluxoscomercio}{487,419}
\newcommand{\Nurfsmaritimas}{103}
\newcommand{\CobSantosPre}{0\%}
\newcommand{\CobSantosPos}{93\%}
\newcommand{\CobRioPre}{27\%}
\newcommand{\CobRioPos}{71\%}
\newcommand{\TdoisBetaFull}{-0.094}
\newcommand{\TdoisPwcbFull}{0.38}
\newcommand{\TquatroBetaFull}{-0.027}
\newcommand{\TquatroPwcbFull}{0.80}
\newcommand{\CabRecTquatroBeta}{-0.188}
\newcommand{\CabRecTquatroPwcb}{0.02}
\newcommand{\CabRecTquatroPct}{17\%}
\newcommand{\PlaceboCabTdoisBeta}{-0.271}
\newcommand{\PlaceboCabTdoisPwcb}{0.01}
\newcommand{\PlaceboCabTquatroPwcb}{0.37}
\newcommand{\FreteBeta}{-0.002}
\newcommand{\FreteP}{0.88}
\newcommand{\SfaPspTl}{0.498}
\newcommand{\SfaPspTlP}{0.02}
\newcommand{\SfaPspPm}{0.206}
\newcommand{\SfaLR}{30.6}
, gerado pelo pipeline (targets::tar_make()) a partir do
% DuckDB. Resultados ainda não estimados aparecem como [RESULT TO BE GENERATED].
%
% Compilação: pdflatex -> bibtex -> pdflatex x2 (ou "Recompile" no Overleaf).
% Opções: [review] = espaço duplo + numeração de linhas (submissão);
%         trocar por [preprint] para a versão de leitura em coluna única.
%==============================================================================
\documentclass[review,authoryear,12pt]{elsarticle}

\usepackage[T1]{fontenc}
\usepackage[utf8]{inputenc}
\usepackage[english]{babel}
\usepackage{lmodern}
\usepackage{microtype}

\usepackage{amsmath,amssymb}
\usepackage{booktabs}
\usepackage{longtable}
\usepackage{graphicx}
\graphicspath{{figures/}}
\usepackage{placeins}
\usepackage{tikz}
\usetikzlibrary{arrows.meta,positioning}
\usepackage{caption}
\captionsetup{font=small,labelfont=bf}
\usepackage[hidelinks]{hyperref}
\usepackage[capitalise,noabbrev]{cleveref}
\crefname{appendix}{Appendix}{Appendices}

\journal{Government Information Quarterly}
\bibliographystyle{elsarticle-harv}

% Macros com números do pipeline (NUNCA editar à mão: é gerado)
% numbers.tex — GERADO AUTOMATICAMENTE por scripts/12_export_overleaf.R.
% Não editar à mão: qualquer alteração é sobrescrita na próxima exportação.
% Gerado em 2026-09-23 12:08:40.72706 a partir de data/processed/observatory.duckdb

\newcommand{\Natracacoes}{1,182,270}
\newcommand{\Anoinicio}{2010}
\newcommand{\Anofim}{2026}
\newcommand{\Npainellinhas}{29,461}
\newcommand{\Nportos}{251}
\newcommand{\Nportospublicos}{34}
\newcommand{\Nportostups}{217}
\newcommand{\Fretemedio}{4.9\%}
\newcommand{\Fretepico}{7.2\%}
\newcommand{\Anofretepico}{2022}
\newcommand{\Nfluxoscomercio}{487,419}
\newcommand{\Nurfsmaritimas}{103}
\newcommand{\CobSantosPre}{0\%}
\newcommand{\CobSantosPos}{93\%}
\newcommand{\CobRioPre}{27\%}
\newcommand{\CobRioPos}{71\%}
\newcommand{\TdoisBetaFull}{-0.094}
\newcommand{\TdoisPwcbFull}{0.38}
\newcommand{\TquatroBetaFull}{-0.027}
\newcommand{\TquatroPwcbFull}{0.80}
\newcommand{\CabRecTquatroBeta}{-0.188}
\newcommand{\CabRecTquatroPwcb}{0.02}
\newcommand{\CabRecTquatroPct}{17\%}
\newcommand{\PlaceboCabTdoisBeta}{-0.271}
\newcommand{\PlaceboCabTdoisPwcb}{0.01}
\newcommand{\PlaceboCabTquatroPwcb}{0.37}
\newcommand{\FreteBeta}{-0.002}
\newcommand{\FreteP}{0.88}
\newcommand{\SfaPspTl}{0.498}
\newcommand{\SfaPspTlP}{0.02}
\newcommand{\SfaPspPm}{0.206}
\newcommand{\SfaLR}{30.6}


% Marcadores visíveis
\newcommand{\TBG}{\textbf{[RESULT TO BE GENERATED]}}
\newcommand{\PROSA}{\textit{[prose to be written by the author from the numbers above]}}

\begin{document}

\begin{frontmatter}

% Título orientado a governo digital / SI (decisão do pesquisador, 11/09: eixo b).
\title{Digitalizing Port Processes: Design and Causal Evaluation\\
of a Port Transformation Observatory}

% Revisão duplamente anônima: SEM \author, \affiliation ou agradecimentos aqui.
% Ver title_page.tex.

\begin{abstract}
% GIQ: máximo de 250 palavras; sem referências; sem abreviaturas não definidas.
% Ideia de abstract em docs/writing_guide.md §2 (usa as macros de numbers.tex).
\TBG
\end{abstract}

% GIQ: 3 a 5 itens, cada um com no máximo 85 caracteres (com espaços).
% Cópia editável em highlights.tex (arquivo separado exigido na submissão).
\begin{highlights}
\item \TBG
\item \TBG
\item \TBG
\end{highlights}

% GIQ: 1 a 7 keywords, em inglês; evitar "and"/"of".
\begin{keyword}
digital government \sep maritime single window \sep information quality \sep
port digitalization \sep staggered difference-in-differences \sep
design science
\end{keyword}

\end{frontmatter}

\section{Introduction}
\label{sec:intro}

% ROTEIRO (protocolo, secao 24.1). A prosa e escrita pelo autor.
% - Tempo em porto como custo de comercio \citep{hummels2013,djankov2010,clark2004}.
% - O porto como sistema interorganizacional de processamento de informacao.
% - O Porto Sem Papel: o DUV substitui ~140 formularios e integra seis anuentes.
% - Variacao identificadora: adocao escalonada 2011--2013 entre portos publicos,
%   com terminais privados (TUPs) nunca tratados no periodo.
% - Contribuicoes: (i) evidencia causal; (ii) artefato de SI (o Observatorio);
%   (iii) medida direta de custo de comercio (frete ad valorem observado).
% - Previa dos achados.

\TBG

\section{Institutional Background}
\label{sec:institutional}

% ROTEIRO: governanca portuaria brasileira (portos organizados vs. TUPs);
% os seis orgaos anuentes de uma escala (autoridade portuaria, Receita Federal,
% ANVISA, MAPA/Vigiagro, Policia Federal, Marinha); a Lei 12.815/2013 como
% choque institucional concomitante, a ser tratado explicitamente.

\TBG

\section{Port Digitalization in Brazil}
\label{sec:digitalization}

% ROTEIRO: taxonomia das intervencoes (janela unica maritima; janela unica de
% comercio exterior; VTMIS; PCS; agendamento; automacao de gate); o que
% distingue transformacao digital de mera informatizacao; o registro de
% intervencoes e sua graduacao de confianca documental.

\input{tables/tab03_intervencoes}

\begin{figure}[htbp]
\centering
\includegraphics[width=\textwidth]{fig10_rollout.pdf}
\caption{Staggered adoption of the Porto Sem Papel single window. Dots mark the date of
the act that made the single window mandatory; bars cover the treated period.
Dashed lines are public ports whose adoption date is not yet documented by an
official act (excluded from estimation). Source: Official Gazette (DOU) and
official news; see \texttt{data/metadata/psp\_portarias\_dou.csv}.}
\label{fig:rollout}
\end{figure}

\FloatBarrier

\section{Economic and Information Systems Theory}
\label{sec:theory}

% ROTEIRO (duas teorias centrais e uma complementar):
% - OIPT \citep{galbraith1974}: a janela unica eleva a capacidade de
%   processamento de informacao e substitui folga (tempo de espera).
% - IS Business Value \citep{melville2004}: capacidades digitais convertem-se em
%   desempenho, condicionadas a recursos complementares \citep{milgrom1990}.
% - TOE \citep{tornatzky1990}: explica adocao e heterogeneidade.
% Hipoteses H1--H10 conforme docs/research_protocol.md.

\TBG

\section{The Brazilian Port Digital Transformation Observatory}
\label{sec:observatory}

% ROTEIRO: o artefato (design science) -- camada de integracao em DuckDB,
% registro curado de intervencoes, auditorias de qualidade, paineis.
% Dimensoes: pessoas, processos, organizacoes, tecnologia.

\TBG

\section{Data}
\label{sec:data}

% ROTEIRO: ANTAQ (Estatistico Aquaviario) e Comex Stat; construcao do painel
% porto--mes; conceitos oficiais de tempo; crosswalk URF--porto e ressalvas.

The operational panel combines \Natracacoes\ vessel calls with cargo
movement recorded by ANTAQ between \Anoinicio\ and \Anofim, aggregated into
\Npainellinhas\ port--month observations covering \Nportos\ ports
(\Nportospublicos\ public ports and \Nportostups\ private terminals).

\input{tables/tab01_amostra}

Operational times follow the official definitions of the agency: $T_1$ is the
interval between arrival at the anchorage and berthing; $T_2$ the wait between
berthing and the start of operations; $T_3$ the cargo operation itself; and
$T_4$ the wait until unberthing, with $T_A = T_2+T_3+T_4$ and $T_E = T_1+T_A$.

\input{tables/tab02_tempos}

\begin{figure}[htbp]
\centering
\includegraphics[width=\textwidth]{fig01_distribuicao_t1.pdf}
\caption{Components of the vessel's stay in port, by year and port type.
Median hours per vessel call with cargo movement (deep-sea, cabotage and
inland navigation). The shaded band marks the single-window rollout
(July 2011--April 2013). Source: ANTAQ waterway statistics, own elaboration.}
\label{fig:t1dist}
\end{figure}

\begin{figure}[htbp]
\centering
\includegraphics[width=\textwidth]{fig11_map.pdf}
\caption{Port facilities in the panel by treatment status. Marker area is proportional to
tonnes handled in 2010--2015. Source: ANTAQ waterway statistics, own elaboration. State boundaries: IBGE (2024).}
\label{fig:map}
\end{figure}

\input{tables/tab04_qualidade}

\FloatBarrier

\section{Measuring Digitalization}
\label{sec:measuring}

% ROTEIRO: regra de datacao (entrada em producao, nunca anuncio); niveis de
% confianca; se e como construir um indice (pesos transparentes, nunca
% ajustados para maximizar significancia).

\TBG

\section{Identification Strategy}
\label{sec:identification}

% ROTEIRO: DiD escalonado robusto a heterogeneidade
% \citep{callaway2021,sunabraham2021}; comparacao not-yet-treated e TUPs como
% nunca tratados (robustez); event study; sensibilidade a pretrends
% \citep{rambachan2023}; synthetic DiD \citep{arkhangelsky2021} para
% intervencoes concentradas; PPML \citep{santossilva2006} para comercio.
% Tabela de estimandos em docs/estimand_table.md.

\TBG

\section{Results}
\label{sec:results}

% EIXO (b), decisao de 11/09: o achado de informacao (H11) abre a secao; os
% efeitos sobre tempos, comercio e eficiencia vem como contexto — tres
% resultados nulos honestos e um efeito documental estreito.

% ATENCAO (23/09): o salto de cobertura e essencialmente Santos (0 -> 67% na janela);
% o efeito medio nao e robusto (DiD A 0,197 p_wcb 0,38; com tendencia por porto ~0).
% Ver tab:pretrend e outputs/tables/coverage_jump_by_port.csv (decisao 42).
\subsection{What the single window changed first: the information the regulator receives}
\label{sec:res_info}

Before the single window, the post-operation waiting time ($T_4$) was
recorded for \CobSantosPre\ of vessel calls in Santos and \CobRioPre\ in Rio
de Janeiro; after it, for \CobSantosPos\ and \CobRioPos\ respectively
(\Cref{fig:coverage}). \Cref{tab:covefeito} reports the corresponding
difference-in-differences estimates.

\input{tables/tab06_cobertura_efeito}

\begin{figure}[htbp]
\centering
\includegraphics[width=\textwidth]{fig13_es_coverage.pdf}
\caption{Event study of the probability that $T_4$ is recorded, dated public ports against
never-treated private terminals, 2010--2015. Sun--Abraham estimator with port and
month fixed effects; month $-1$ is the reference. Shaded band: 95\% uniform
(sup-$t$) band \citep{montielolea2019}; vertical lines: pointwise 95\% intervals; standard errors
clustered by port. Source: ANTAQ waterway statistics, own elaboration.}
\label{fig:escoverage}
\end{figure}

\PROSA

\subsection{Operational times: no aggregate effect}
\label{sec:res_times}

Across dated public ports, the static effect of the single window on
documentary waiting times is \TdoisBetaFull\ for $T_2$ (wild-bootstrap
$p=\TdoisPwcbFull$) and \TquatroBetaFull\ for $T_4$ ($p=\TquatroPwcbFull$);
the physical queue ($T_1$) shows no movement (\Cref{tab:inferencia}).

\input{tables/tab07_inferencia}

\begin{figure}[htbp]
\centering
\includegraphics[width=\textwidth]{fig02_coortes_psp.pdf}
\caption{Documentary waiting times by adoption cohort, quarterly medians,
2010--2015. Dashed lines mark the first production date of each cohort.
Raw medians mix changes in times with changes in which calls are recorded
(\Cref{sec:res_info}); descriptive evidence only. Source: ANTAQ waterway statistics, own elaboration.}
\label{fig:cohorts}
\end{figure}

\PROSA

% ATENCAO (23/09, decisao do autor pendente): o titulo e o paragrafo abaixo
% foram escritos antes do teste de pre-tendencias. Com tendencia linear por porto
% e com HonestDiD (M=0) o efeito da cabotagem NAO sobrevive (tab:pretrend,
% fig:honestdid, docs/decisions_log.md decisao 42). Reescrever ao redigir a prosa.
\subsection{Where a documentary effect survives: recurrent cabotage calls}
\label{sec:res_nav}

The only estimate that survives both few-cluster inference and the
anticipation placebo is the post-operation waiting time of recurrent
cabotage vessels: \CabRecTquatroBeta\ in logs (about \CabRecTquatroPct\
shorter, wild-bootstrap $p=\CabRecTquatroPwcb$; \Cref{tab:heterog}).

\input{tables/tab08_heterogeneidade}

\begin{figure}[htbp]
\centering
\includegraphics[width=\textwidth]{fig12_es_cabotage.pdf}
\caption{Event study of documentary waiting times for recurrent cabotage vessels (three or
more calls at the same port in the previous twelve months), design A
(not-yet-treated public ports, 2010--2013). Sun--Abraham estimator with port,
month and navigation fixed effects and cargo controls; month $-1$ is the
reference. Shaded band: 95\% uniform (sup-$t$) band \citep{montielolea2019}; vertical lines: pointwise
95\% intervals; standard errors clustered by port. Source: ANTAQ waterway statistics, own elaboration.}
\label{fig:escabotage}
\end{figure}

\PROSA

\subsection{Long-run dynamics, 2010--2026}
\label{sec:res_long}

\begin{figure}[htbp]
\centering
\includegraphics[width=\textwidth]{fig08_horizonte_longo.pdf}
\caption{Long-horizon event study for documentary waiting times, dated public
ports against never-treated private terminals, 2010--2026. Sun--Abraham
estimator with port, month and navigation fixed effects and cargo controls.
Shaded band: 95\% uniform (sup-$t$) band \citep{montielolea2019}; vertical lines: pointwise 95\%
intervals; standard errors clustered by port. Source: ANTAQ waterway statistics, own elaboration.}
\label{fig:longrun}
\end{figure}

\PROSA

\FloatBarrier

\section{Mechanisms}
\label{sec:mechanisms}

% ROTEIRO: sequencia digitalizacao -> tempo/incerteza -> eficiencia ->
% movimentacao/comercio, testada com o mesmo desenho. Sem alegacao de
% mediacao causal formal.

\TBG

\section{Heterogeneity}
\label{sec:heterogeneity}

% ROTEIRO: publico vs. privado; conteiner vs. granel; porte; regiao;
% infraestrutura; numero de orgaos integrados; qualidade da implementacao.

\TBG

\section{Robustness}
\label{sec:robustness}

% Cada teste responde a uma ameaca nomeada. WCB e permutacao: poucos clusters
% (tab07, na secao de resultados). Placebo de antecipacao: datacao e
% pre-tendencias. Honest DiD: violacoes limitadas de tendencias paralelas.

The anticipation placebo shifts each treatment date twelve months earlier
and is estimated on the true pre-period only. For cabotage $T_2$ it returns
\PlaceboCabTdoisBeta\ ($p=\PlaceboCabTdoisPwcb$): the apparent reduction in
$T_2$ predates adoption and is therefore a pre-trend, not an effect. The
$T_4$ placebo is not distinguishable from zero ($p=\PlaceboCabTquatroPwcb$).

\input{tables/tab09_placebo}

\input{tables/tab12_lee_bounds}

\begin{figure}[htbp]
\centering
\includegraphics[width=\textwidth]{fig14_estimates.pdf}
\caption{Static difference-in-differences estimates by sample (design A). Dots and thin
lines: point estimate and 95\% cluster-robust interval; wide shaded bars: Lee
bounds from \Cref{tab:lee} where computed; $p$ (right): wild cluster
bootstrap $p$-value (Webb weights). Source: ANTAQ waterway statistics, own elaboration.}
\label{fig:estimates}
\end{figure}

\input{tables/tab13_pretrend}

\begin{figure}[htbp]
\centering
\includegraphics[width=\textwidth]{fig15_honestdid.pdf}
\caption{Robust confidence intervals for the average cabotage effect (event quarters 0--4)
under the smoothness restriction of \citet{rambachan2023}. $M=0$ allows the pre-adoption
trend to continue linearly after adoption; larger $M$ allows curvature. ``Original'' is
the conventional 95\% interval. Design A, Sun--Abraham event study in quarters with
port, month and navigation fixed effects and cargo controls; standard errors clustered
by port. Source: ANTAQ waterway statistics, own elaboration.}
\label{fig:honestdid}
\end{figure}

\PROSA

\PROSA

\FloatBarrier

\section{Information Systems Contributions}
\label{sec:is}

% ROTEIRO: artefato e avaliacao causal do fenomeno que ele mede; implicacoes
% para OIPT e IS business value em contextos interorganizacionais de governo.

\TBG

\subsection{The single window and the quality of official statistics}
\label{sec:infoquality}

% ROTEIRO (H11, decisao do pesquisador em 2026-09-10): a decomposicao
% documental dos tempos (T2/T4) passa a ser reportada a ANTAQ junto com a
% entrada do Porto Sem Papel no maior porto do pais; em outros portos a
% cobertura muda em momentos nao relacionados ao tratamento. Implicacoes:
% (i) o artefato digital altera o processo gerador das estatisticas que o
% avaliam (information quality); (ii) consequencia metodologica -- restricao
% da analise de T2/T4 aos portos de cobertura estavel.

Coverage of the documentary time decomposition is far from uniform.
In Santos, the share of vessel calls with a recorded $T_4$ rises from
\CobSantosPre\ in 2010--2011 to \CobSantosPos\ from 2012 onwards, immediately
after the single window entered production; in Rio de Janeiro it rises from
\CobRioPre\ to \CobRioPos\ in 2014. Elsewhere, coverage changes are unrelated
to treatment timing (\Cref{fig:coverage}).

\begin{figure}[htbp]
\centering
\includegraphics[width=\textwidth]{fig07_cobertura_t4.pdf}
\caption{Share of vessel calls with a recorded post-operation waiting time
($T_4$), by port and year. Year of Porto Sem Papel production start in
parentheses.}
\label{fig:coverage}
\end{figure}

\TBG

\FloatBarrier

\section{Implications for digital government}
\label{sec:implications}

% Enquadramento GIQ (decisão 40): o que o caso diz para políticas de governo
% digital — (i) uma janela única muda primeiro o REGISTRO (cobertura de T4)
% e só depois, e de forma localizada, o processo (cabotagem recorrente);
% (ii) avaliações que ignoram a mudança de medição atribuem à digitalização
% efeitos que são artefatos; (iii) a promessa de custos de comércio menores
% não se verifica no horizonte observado (\Cref{app:trade}).
% Fronteira e comércio detalhados nos \Cref{app:efficiency,app:trade}.

\PROSA

\FloatBarrier

\section{Limitations}
\label{sec:limitations}

% ROTEIRO: pre-periodo curto da coorte de 2011; todos os portos publicos
% tratados ate 2013 (identificacao de longo prazo limitada); URF nao e porto;
% qualidade do reporte dos tempos; decomposicao T2--T4 incompleta nos
% primeiros anos.

\TBG

\input{sections/17_conclusion}

% Declarações exigidas pela GIQ ficam na folha de rosto (title_page.tex):
% competing interests, funding, CRediT, generative AI, data availability.

\clearpage
% RASCUNHO: lista todas as entradas do .bib para o autor ver o que está
% disponível. REMOVER \nocite{*} na versão de submissão.
\nocite{*}
\bibliography{refs}

\clearpage
\appendix
\section{Appendix}
\label{sec:appendix}

\subsection{Estimand table}
% Portar de docs/estimand_table.md quando congelada.
\TBG

\subsection{Intervention registry}
% Registro completo com fontes e graus de confianca.
\TBG

\subsection{Additional descriptive evidence}

\begin{figure}[htbp]
\centering
\includegraphics[width=\textwidth]{fig04_movimentacao.pdf}
\caption{Annual cargo throughput recorded in the panel, by port type.}
\label{fig:throughput}
\end{figure}

\FloatBarrier

\section{Stochastic frontier model}
\label{app:efficiency}

% Apêndice B (GIQ: apêndices A, B, C; equações como Eq. (B.1)). Movido do
% corpo do texto em 12/09 (decisão 40): na GIQ a fronteira é evidência
% complementar sobre eficiência técnica, não o argumento central.
% Fronteira estocástica \citep{aigner1977,battese1995}: output ln toneladas;
% inputs ln horas de berço, ln berços ativos; ineficiência meio-normal com
% variância heteroscedástica (PSP, share de contêiner). Ver docs/efficiency_model.md.

Public ports only (34 ports), 2010--2025. The translog specification is
preferred (LR $=\SfaLR$). The single window enters the inefficiency variance
with coefficient \SfaPspTl\ ($p=\SfaPspTlP$) in the port-year translog and
\SfaPspPm\ in the port-month specification (\Cref{tab:sfa}). Draft and berth
capacity are not observed; the estimates are associations, not causal effects
(see \Cref{sec:limitations}).

\input{tables/tab11_fronteira}

\PROSA

\FloatBarrier

\section{Trade values and observed freight costs}
\label{app:trade}

% Apêndice C. Movido da §15 em 12/09 (decisão 40): resultado nulo sobre
% custos de comércio — relevante para a GIQ como limite da promessa da
% janela única, mas não é o núcleo do argumento.

Import freight charges provide a directly observed measure of trade costs.
Across maritime imports, declared freight averaged \Fretemedio\ of FOB value
over the sample period, peaking at \Fretepico\ in \Anofretepico.

\begin{figure}[htbp]
\centering
\includegraphics[width=\textwidth]{fig03_frete_advalorem.pdf}
\caption{Ad valorem freight rate on Brazilian maritime imports.}
\label{fig:freight}
\end{figure}

The single window has no detectable effect on trade values or on the
observed freight rate (\Cref{tab:comercio}): the coefficient on the log
freight-to-FOB ratio is \FreteBeta\ ($p=\FreteP$).

\input{tables/tab10_comercio}

% Emissões de fundeio: SOMENTE com fatores citáveis (IMO GHG Study); como o
% efeito sobre T1 é nulo, não há espera evitada a valorar.

\PROSA

\FloatBarrier


\end{document}
